\documentclass[a4paper]{article}
\usepackage{amsfonts}
\usepackage{amsmath}
\usepackage{amssymb}
\usepackage{graphicx}     

\begin{document}
  
\noindent \large Auteur: Abdellah El Moussaoui \\
\noindent \large Studierichting: Informatica\\
\noindent \large Oefeningenreeks 1C nr. 45.10\\

\medskip

\normalsize

\medskip

Trek de derdemachtswortels in $\mathbb{C}$:\\

$z^3=4\sqrt{2}(1+i)$\\

Voor $1+i$ hebben we:\\

$r=\sqrt{1^2+1^2}=\sqrt{2}$\\

$\theta=\operatorname{Bgtan}\left(\dfrac{1}{1}\right)=\dfrac{\pi}{4}$\\

$\Rightarrow 1+i=\sqrt{2}\operatorname{cis}\left(\dfrac{\pi}{4}\right)$\\

$\Rightarrow z^3=4\sqrt{2}\cdot \sqrt{2}\operatorname{cis}\left(\dfrac{\pi}{4}\right)=8\operatorname{cis}\left(\dfrac{\pi}{4}\right)$\\

$\Rightarrow z_k=2\operatorname{cis}\left(\dfrac{\dfrac{\pi}{4}+2k\pi}{3}\right),\ k=0,1,2$\\

$\Rightarrow z_k=2\operatorname{cis}\left(\dfrac{\pi}{12}+\dfrac{2k\pi}{3}\right),\ k=0,1,2$\\

$\Rightarrow z_0=2\operatorname{cis}\left(\dfrac{\pi}{12}\right)=\dfrac{\sqrt6+\sqrt{2}}{2}+i\dfrac{\sqrt6-\sqrt{2}}{2}$\\

$\Rightarrow z_1=2\operatorname{cis}\left(\dfrac{3\pi}{4}\right)=-\sqrt{2}+i\sqrt{2}$\\

$\Rightarrow z_2=2\operatorname{cis}\left(\dfrac{17\pi}{12}\right)=-\dfrac{\sqrt6-\sqrt{2}}{2}-i\dfrac{\sqrt6+\sqrt{2}}{2}$\\

$\Rightarrow \boxed{z_0=\dfrac{\sqrt6+\sqrt{2}}{2}+i\dfrac{\sqrt6-\sqrt{2}}{2},\ z_1=-\sqrt{2}+i\sqrt{2},\ z_2=-\dfrac{\sqrt6-\sqrt{2}}{2}-i\dfrac{\sqrt6+\sqrt{2}}{2}}$\\


\begin{thebibliography}{999}
%\bibitem{a} Referentie
\end{thebibliography}
\end{document} 